\begin{ex}
Tìm số tròn chục lớn nhất có hai chữ số.
\shortans{90}
\loigiai{
Số tròn chục lớn nhất có hai chữ số là số $90$.
\textbf{Đáp số:} $90$.
}
\end{ex}

\begin{ex}
Tìm số lớn nhất có hai chữ số.
\shortans{99}
\loigiai{
Số lớn nhất có hai chữ số là số $99$.
\textbf{Đáp số:} $99$.
}
\end{ex}

\begin{ex}
Tìm số liền sau của số lớn nhất có hai chữ số.
\shortans{100}
\loigiai{
Số liền sau của 99 là số $100$.
\textbf{Đáp số:} $100$.
}
\end{ex}

\begin{ex}
Số 75 gồm mấy chục và mấy đơn vị? Lấy chữ số hàng chục trừ đi chữ số hàng đơn vị được kết quả là bao nhiêu?
\shortans{2}
\loigiai{
Số 75 có chữ số hàng chục là 7, hàng đơn vị là 5: $7 - 5 = 2$.
\textbf{Đáp số:} $2$.
}
\end{ex}

\begin{ex}
Cho ba chữ số: 3, 0, 8. Lập số lớn nhất có hai chữ số từ các chữ số đã cho.
\shortans{83}
\loigiai{
Chọn hàng chục lớn nhất là 8, hàng đơn vị lớn nhất là 3: số đó là $83$.
\textbf{Đáp số:} $83$.
}
\end{ex}

\begin{ex}
Cho ba chữ số: 3, 0, 8. Lập số bé nhất có hai chữ số khác nhau từ các chữ số đã cho.
\shortans{30}
\loigiai{
Hàng chục bé nhất khác 0 là 3, hàng đơn vị bé nhất là 0: số đó là $30$.
\textbf{Đáp số:} $30$.
}
\end{ex}

\begin{ex}
Cho các số: $56, 65, 59, 62$. Sắp xếp theo thứ tự từ bé đến lớn, số đứng ở vị trí thứ ba là số nào?
\shortans{62}
\loigiai{
Thứ tự từ bé đến lớn: $56 < 59 < 62 < 65$. Số đứng thứ ba là số $62$.
\textbf{Đáp số:} $62$.
}
\end{ex}

\begin{ex}
Tìm số tự nhiên $a$ biết rằng: $69 < a < 71$.
\shortans{70}
\loigiai{
Số tự nhiên nằm giữa 69 và 71 là số tròn chục $70$.
\textbf{Đáp số:} $70$.
}
\end{ex}

\begin{ex}
Bác nông dân có 5 chục quả trứng gà và 3 chục quả trứng vịt. Hỏi bác nông dân có tất cả bao nhiêu quả trứng cả hai loại?
\shortans{80}
\loigiai{
Tổng số chục trứng là: $5 + 3 = 8$ (chục).\n8 chục quả trứng tương ứng với $80$ quả trứng.
\textbf{Đáp số:} $80$.
}
\end{ex}

\begin{ex}
Một cửa hàng có 90 chiếc cặp sách, đã bán được 40 chiếc cặp sách. Hỏi cửa hàng còn lại bao nhiêu chiếc cặp sách?
\shortans{50}
\loigiai{
Số cặp sách còn lại là: $90 - 40 = 50$ (chiếc cặp sách).
\textbf{Đáp số:} $50$.
}
\end{ex}

\begin{ex}
Tìm số tròn chục $x$ thỏa mãn: $20 + x = 70$. Giá trị của $x$ là bao nhiêu?
\shortans{50}
\loigiai{
Ta có: $x = 70 - 20 = 50$.
\textbf{Đáp số:} $50$.
}
\end{ex}

\begin{ex}
Tìm số tròn chục $y$ thỏa mãn: $80 - y = 30$. Giá trị của $y$ là bao nhiêu?
\shortans{50}
\loigiai{
Ta có: $y = 80 - 30 = 50$.
\textbf{Đáp số:} $50$.
}
\end{ex}

\begin{ex}
Bạn Lan có 35 que tính, bạn Mai có 38 que tính. Hỏi bạn Mai có nhiều hơn bạn Lan bao nhiêu que tính?
\shortans{3}
\loigiai{
Số que tính Mai nhiều hơn Lan là: $38 - 35 = 3$ (que tính).
\textbf{Đáp số:} $3$.
}
\end{ex}

\begin{ex}
Lớp 1A có 32 học sinh, lớp 1B có 35 học sinh. Hỏi cả hai lớp có tất cả bao nhiêu học sinh?
\shortans{67}
\loigiai{
Tổng số học sinh cả hai lớp là: $32 + 35 = 67$ (học sinh).
\textbf{Đáp số:} $67$.
}
\end{ex}

\begin{ex}
Tìm số có hai chữ số giống nhau biết tổng hai chữ số bằng 8.
\shortans{44}
\loigiai{
Hai chữ số giống nhau có tổng bằng 8 là 4 và 4. Số đó là $44$.
\textbf{Đáp số:} $44$.
}
\end{ex}

\begin{ex}
Tìm số có hai chữ số giống nhau biết tổng hai chữ số bằng 12.
\shortans{66}
\loigiai{
Hai chữ số giống nhau có tổng bằng 12 là 6 và 6. Số đó là $66$.
\textbf{Đáp số:} $66$.
}
\end{ex}

\begin{ex}
Hiệu của số tròn chục lớn nhất và số tròn chục nhỏ nhất có hai chữ số là bao nhiêu?
\shortans{80}
\loigiai{
Số tròn chục lớn nhất là 90, nhỏ nhất là 10: $90 - 10 = 80$.
\textbf{Đáp số:} $80$.
}
\end{ex}

\begin{ex}
Từ số 40 đến số 50 có tất cả bao nhiêu số tự nhiên liên tiếp?
\shortans{11}
\loigiai{
Số lượng số là: $50 - 40 + 1 = 11$ (số).
\textbf{Đáp số:} $11$.
}
\end{ex}

\begin{ex}
Một người có 60 bông hoa, người đó cắm vào các lọ mỗi lọ 10 bông hoa. Hỏi cắm được tất cả bao nhiêu lọ hoa?
\shortans{6}
\loigiai{
60 bông hoa gồm 6 chục bông hoa, chia mỗi lọ 1 chục nên cắm được $6$ lọ hoa.
\textbf{Đáp số:} $6$.
}
\end{ex}

\begin{ex}
Đoàn tàu hỏa có 10 toa xe. Toa thứ 7 tính từ đầu tàu là toa thứ mấy tính từ đuôi tàu lên?
\shortans{4}
\loigiai{
Vị trí từ đuôi tàu lên là: $10 - 7 + 1 = 4$.
\textbf{Đáp số:} $4$.
}
\end{ex}

\begin{ex}
Tìm số tự nhiên $m$ có hai chữ số thỏa mãn: chữ số hàng chục hơn chữ số hàng đơn vị là 7 và hàng đơn vị là 1.
\shortans{81}
\loigiai{
Chữ số hàng chục là: $1 + 7 = 8$. Số đó là $81$.
\textbf{Đáp số:} $81$.
}
\end{ex}

\begin{ex}
Tính tổng của 4 chục và 5 chục.
\shortans{90}
\loigiai{
Ta có: 4 chục + 5 chục = 9 chục, tức là $90$.
\textbf{Đáp số:} $90$.
}
\end{ex}

\begin{ex}
Tính hiệu của 9 chục và 6 chục.
\shortans{30}
\loigiai{
Ta có: 9 chục - 6 chục = 3 chục, tức là $30$.
\textbf{Đáp số:} $30$.
}
\end{ex}

\begin{ex}
Số nào cộng với 50 thì bằng 90?
\shortans{40}
\loigiai{
Số đó là: $90 - 50 = 40$.
\textbf{Đáp số:} $40$.
}
\end{ex}

\begin{ex}
Số nào trừ đi 20 thì bằng 60?
\shortans{80}
\loigiai{
Số đó là: $60 + 20 = 80$.
\textbf{Đáp số:} $80$.
}
\end{ex}
